% Étui de décontamination des lentilles avec disque catalytique (exercice 14, chapitre 4)
\begin{tikzpicture}[x=1cm,y=1cm]
% liquide
\fill[pcEau] (-0.74,2.55) -- (-0.74,0.35) arc[start angle=180,end angle=360,radius=0.74] -- (0.74,2.55) -- cycle;
\draw[pcSurface] (-0.74,2.55) -- (0.74,2.55);
% tige et paniers porte-lentilles
\fill[otInk!35] (-0.05,0.95) rectangle (0.05,3.05);
\foreach \s in {-1,1}{%
  \draw[otNavy!70,line width=0.5pt,dash pattern=on 1.4pt off 0.9pt] (\s*0.06,2.35) .. controls (\s*0.62,2.32) and (\s*0.62,1.38) .. (\s*0.06,1.35);
  \draw[otGreenDark,line width=0.9pt] (\s*0.12,2.2) .. controls (\s*0.45,2.1) and (\s*0.45,1.6) .. (\s*0.12,1.5);}
% disque catalytique
\filldraw[fill=otInk!65,draw=otInk!85,line width=0.4pt,rounded corners=0.3mm] (-0.5,0.78) rectangle (0.5,0.95);
% bulles de dioxygène
\foreach \bx/\by/\br in {-0.38/1.1/0.05,0.3/1.18/0.045,-0.12/1.42/0.04,0.5/1.6/0.05,-0.55/1.85/0.045,0.22/2.05/0.04,-0.3/2.3/0.05,0.45/2.38/0.04}{%
  \draw[otNavy!60,line width=0.4pt,fill=white] (\bx,\by) circle (\br);}
% récipient et couvercle
\draw[pcVerre] (-0.8,3.05) -- (-0.8,0.35) arc[start angle=180,end angle=360,radius=0.8] -- (0.8,3.05);
\filldraw[pcAppareil,rounded corners=0.6mm] (-0.95,3.05) rectangle (0.95,3.5);
\fill[white] (0.42,3.38) rectangle (0.52,3.5);
% légendes
\pcNote[west]{(0.5,3.45)}{(1.5,3.75)}{couvercle muni d'un évent}
\pcNote[west]{(0.38,1.85)}{(1.5,2.05)}{lentilles dans leurs paniers}
\pcNote[west]{(0.74,2.55)}{(1.5,2.75)}{solution de peroxyde d'hydrogène à 3 \%}
\pcNote[west]{(0.5,0.86)}{(1.5,0.86)}{disque recouvert de platine}
\pcNote[east]{(-0.55,1.85)}{(-1.4,1.85)}{bulles de dioxygène}
\end{tikzpicture}
